% =====================================================================================
% Shared style for the three papers of the Trivijaya Quant Research programme.
%
% The three papers report one programme and must look like it: same typography, same
% heading hierarchy, same table rules, same panel vocabulary, same cover structure.
% Anything paper-specific (title, running head, abstract) stays in the paper's own file.
%
% Each paper defines, BEFORE inputting this file:
%   \paperName        e.g. AlphaAudit
%   \paperSubtitle    the one-line question the paper answers
%   \paperNumber      1, 2 or 3
%   \paperShortHead   the right-hand running head
% =====================================================================================

\usepackage[utf8]{inputenc}
\usepackage[T1]{fontenc}
\usepackage{lmodern}
\usepackage[scaled=0.92]{helvet}          % sans face, headings only
\usepackage[margin=1.05in,top=1.15in,bottom=1.1in]{geometry}
\usepackage{microtype}
\usepackage{amsmath,amssymb}
\usepackage{booktabs}
\usepackage{array}
\usepackage{graphicx}
\usepackage{xcolor}
\usepackage[most]{tcolorbox}
\usepackage{titlesec}
\usepackage{enumitem}
\usepackage[font=small,labelfont={bf,color=navy},labelsep=period]{caption}
\usepackage{fancyhdr}
\usepackage{natbib}
\usepackage[colorlinks=true]{hyperref}   % loaded last, as hyperref requires

% The rupee sign is not in the default T1 encoding. A plain "Rs." is unambiguous, renders
% identically on every machine, and avoids a font dependency a reader would have to satisfy
% before they could compile the paper that asks them to check our numbers.
\providecommand{\rupee}{\textrm{Rs.}\,}

% ---- Palette: white ground, blues, one warm accent for our own errors ---------------
\definecolor{navy}{HTML}{0E2E5C}
\definecolor{steel}{HTML}{2C6BA8}
\definecolor{sky}{HTML}{5B93C7}
\definecolor{mist}{HTML}{F0F5FB}
\definecolor{ash}{HTML}{F6F7F9}
\definecolor{ink}{HTML}{1A1A1A}
\definecolor{rust}{HTML}{9C4221}

\color{ink}
\hypersetup{linkcolor=steel, citecolor=steel, urlcolor=steel}

% ---- Repository, cited often enough to deserve one definition ----------------------
\newcommand{\repourl}{https://github.com/Ashwin-aadi/Trivijaya-Quant-Research}
\newcommand{\repo}{\href{\repourl}{\texttt{github.com/Ashwin-aadi/Trivijaya-Quant-Research}}}
\newcommand{\siteurl}{https://ashwin-aadi.github.io/Trivijaya-Quant-Research/}
\newcommand{\site}{\href{\siteurl}{\texttt{ashwin-aadi.github.io/Trivijaya-Quant-Research}}}

% ---- Headings ----------------------------------------------------------------------
\newcommand{\sectionrule}{{\color{sky}\rule{\linewidth}{0.5pt}}\vspace{-0.9em}}
\titleformat{\section}[block]
  {\sffamily\large\bfseries\color{navy}}
  {\colorbox{navy}{\textcolor{white}{\sffamily\bfseries\ \thesection\ }}}{0.7em}{}
  [\vspace{0.25em}\sectionrule]
\titleformat{\subsection}[block]
  {\sffamily\normalsize\bfseries\color{navy}}{\thesubsection}{0.6em}{}
\titleformat{\paragraph}[runin]
  {\sffamily\bfseries\color{navy}}{}{0em}{}[\;]
\titlespacing*{\section}{0pt}{1.6em}{0.8em}
\titlespacing*{\subsection}{0pt}{1.2em}{0.5em}

% ---- Running head ------------------------------------------------------------------
\pagestyle{fancy}\fancyhf{}
\renewcommand{\headrulewidth}{0.4pt}
\renewcommand{\headrule}{{\color{sky}\hrule height 0.4pt}}
\fancyhead[L]{\sffamily\scriptsize\color{navy}\MakeUppercase{\paperName}}
\fancyhead[R]{\sffamily\scriptsize\color{navy}\paperShortHead}
\fancyfoot[C]{\sffamily\scriptsize\color{navy}\thepage}
\fancypagestyle{plain}{\fancyhf{}\renewcommand{\headrulewidth}{0pt}
  \fancyfoot[C]{\sffamily\scriptsize\color{navy}\thepage}}

% ---- Panels ------------------------------------------------------------------------
% A finding the reader must not skim past.
\newtcolorbox{finding}[1][]{
  enhanced, breakable, colback=mist, colframe=steel, boxrule=0pt,
  leftrule=2.5pt, arc=1pt, left=9pt, right=9pt, top=7pt, bottom=7pt,
  coltext=navy, fonttitle=\sffamily\bfseries\small, #1}

% A place where we report our own error. Deliberately a different colour from the findings,
% so a reader flipping through can see at a glance that the paper contains both.
\newtcolorbox{selfreport}[1][]{
  enhanced, breakable, colback=ash, colframe=rust, boxrule=0pt,
  leftrule=2.5pt, arc=1pt, left=9pt, right=9pt, top=7pt, bottom=7pt,
  fonttitle=\sffamily\bfseries\small, #1}

% What a result does *not* license. Every headline in this programme has one of these, because
% the distance between "we measured X on our corpus" and "X is true of AI" is where this
% literature goes wrong.
\newtcolorbox{scopebox}[1][]{
  enhanced, breakable, colback=white, colframe=sky, boxrule=0.6pt,
  arc=1pt, left=9pt, right=9pt, top=6pt, bottom=6pt,
  fonttitle=\sffamily\bfseries\small, title={\sffamily\bfseries What this does not say}, #1}

\newcommand{\keyfinding}[1]{\begin{finding}#1\end{finding}}

% A one-line takeaway that sits directly against a figure or table.
\newcommand{\takeaway}[1]{%
  \vspace{-0.4em}\noindent{\small\sffamily\color{navy}\textbf{Takeaway.}\ #1}\vspace{0.5em}}

% ---- Tables ------------------------------------------------------------------------
\newcommand{\thead}[1]{{\sffamily\bfseries\color{navy}#1}}
\renewcommand{\arraystretch}{1.15}

\setlist[itemize]{leftmargin=1.3em, itemsep=0.25em, topsep=0.4em}
\setlist[enumerate]{leftmargin=1.6em, itemsep=0.35em, topsep=0.4em}

% ---- Evidence-scope tags -----------------------------------------------------------
% The programme's central discipline: every finding is tagged with the population it was
% measured on, so a local-corpus observation can never be read as a statement about AI.
\newcommand{\tagbox}[2]{{\sffamily\scriptsize\textbf{\textcolor{#1}{[#2]}}}}
\newcommand{\taglocal}{\tagbox{steel}{LOCAL 7B}}
\newcommand{\tagfrontier}{\tagbox{navy}{FRONTIER}}
\newcommand{\tagmethod}{\tagbox{navy}{GENERATION METHOD}}
\newcommand{\tagindep}{\tagbox{rust}{GENERATOR-INDEPENDENT}}
\newcommand{\tagexpl}{\tagbox{rust}{EXPLORATORY}}
\newcommand{\tagprereg}{\tagbox{steel}{PRE-REGISTERED}}

% ---- Cover page --------------------------------------------------------------------
\newcommand{\programmecover}[1]{%
\begin{titlepage}
\thispagestyle{empty}
\raggedright
\vspace*{0.6in}
{\sffamily\small\color{steel}\MakeUppercase{Trivijaya Quant Research Programme}\\[0.2em]
 \color{sky}\rule{\linewidth}{1.2pt}}\\[0.3em]
{\sffamily\small\color{steel}Paper \paperNumber\ of 3 \quad$\cdot$\quad
 Measurement infrastructure for AI-driven quantitative research on Indian equities}

\vspace{1.5in}
{\sffamily\bfseries\color{navy}\fontsize{46}{50}\selectfont \paperName}\\[0.7em]
{\color{sky}\rule{0.4\linewidth}{2pt}}\\[0.9em]
{\sffamily\Large\color{navy} \paperSubtitle}

\vspace{0.9in}
{\sffamily\large\color{ink} Ashwin Jain}\\[0.35em]
{\rmfamily\small Department of Computer Science and Engineering\\
Motilal Nehru National Institute of Technology Allahabad\\
Prayagraj, India\\[0.2em]
\texttt{ashwinaadi2007@gmail.com}}

\vfill
{\small\sffamily\color{navy}#1}

\vspace{0.5em}
{\small\rmfamily Code, corpus, fixtures and evaluation protocol: \repo\\
Programme overview and results: \site}\\[0.6em]
{\sffamily\small\color{steel} 8 August 2026}
\end{titlepage}}

% ---- The trilogy map, printed identically in all three papers ----------------------
\newcommand{\trilogymap}{%
\begin{tcolorbox}[enhanced, colback=ash, colframe=sky, boxrule=0.6pt, arc=1pt,
  left=10pt, right=10pt, top=8pt, bottom=8pt]
\small\sffamily\color{navy}
\textbf{One programme, three questions.} Each paper answers the question the previous one
raises, on one shared universe, one shared cost model and one shared backtest engine.
\vspace{0.5em}

\begin{tabular}{@{}p{0.055\linewidth}p{0.22\linewidth}p{0.65\linewidth}@{}}
\textbf{1} & \textbf{AlphaAudit} & Can an AI produce a trading strategy that is real,
  rather than a plausible artefact of luck and leakage? \\[0.3em]
\textbf{2} & \textbf{RegimeStress} & If a strategy looks real, when does it break? \\[0.3em]
\textbf{3} & \textbf{FlowState} & If a strategy survives, how much capital can it absorb? \\
\end{tabular}
\vspace{0.4em}

\textbf{Code $\rightarrow$ validity $\rightarrow$ robustness $\rightarrow$ scale.} The
programme's overarching question is whether making the AI better makes the \emph{research}
better --- and the three papers answer it only to the extent their experiments support.
\end{tcolorbox}}
